\documentclass[11pt]{paper}

\usepackage[T1]{fontenc}

\usepackage{graphicx}
\usepackage{amssymb}
\usepackage{amstext}
\usepackage{amsmath}
\usepackage{a4wide,color}
\usepackage[utf8]{inputenc}
\usepackage[frenchb]{babel}
\usepackage{xspace}
\usepackage{anysize}
\usepackage{tabularx}
\usepackage{multirow}
\usepackage{fancybox}
\usepackage{fancyhdr}
\usepackage{bbding}

\usepackage{color}
\usepackage{float}

\usepackage[toc,page]{appendix} 
\usepackage{lscape}

\usepackage{placeins}

\usepackage{listings}

\newcommand{\version}{1.0.0{ }}

\lstdefinestyle{customc}{
  belowcaptionskip=1\baselineskip,
  breaklines=true,
  frame=L,
  xleftmargin=\parindent,
  language=C,
  showstringspaces=false,
  basicstyle=\tiny\ttfamily,
  keywordstyle=\bfseries\color{green!40!black},
  commentstyle=\itshape\color{purple!40!black},
  identifierstyle=\color{blue},
  stringstyle=\color{orange},
  tabsize=2
}

\usepackage{todonotes}

\usepackage[framemethod=TikZ]{mdframed}

\lstset{% general command to set parameter(s)
	basicstyle=\footnotesize, % print whole listing small
	keywordstyle=\color{magenta}\bfseries, % underlined bold black keywords
	identifierstyle=, % nothing happens
	commentstyle=\color{black}, % white comments
	stringstyle=\ttfamily, % typewriter type for strings
	showstringspaces=false % no special string spaces
} 

\lstdefinelanguage{aald}
	{morekeywords={system, implementation},
	sensitive=false,
	morecomment=[l]{//},
	morecomment=[s]{/*}{*/},
	morestring=[b]",
}

\usepackage[colorlinks=true]{hyperref}

\usepackage[normalem]{ulem}
\usepackage{color}

\definecolor{Fond}{gray}{0.7}

\renewcommand{\floatpagefraction}{.99}
\renewcommand{\textfraction}{.01}
\newcommand{\modif}[1]{\textcolor{red}{\uline{#1}}}


\pagestyle{fancy}
\fancyhf{}
\fancyhead[RE,RO]{\thepage}
\fancyhead[LE]{}
\fancyhead[LO]{Temps réel }
\fancyfoot[LO]{Centrale Nantes -- \today}
\fancypagestyle{plain}{%
  \fancyhf{} % get rid of headers
  \renewcommand{\headrulewidth}{0pt} % and the line
}

\newenvironment{maliste}%
{ \begin{list}%
	{\ArrowBoldRightStrobe}%
	{\setlength{\labelwidth}{30pt}%
	 \setlength{\leftmargin}{35pt}%
	 \setlength{\itemsep}{\parsep}}}%
{ \end{list} }

 \setlength{\parskip}{1.2ex plus .9ex minus .4ex}

\newcommand{\ocv}{{\tt OpenCV}\xspace} 
\newcommand{\alsa}{{\tt ALSA}\xspace} 
\newcommand{\aadl}{{\tt AADL}\xspace} 
\newcommand{\xeno}{{\tt Xenomai}\xspace} 


\title{{\Large Systèmes interconnectés -- SINT }\\
\vspace{-0.4em}
{\large Sujet de TP 0 : Prise en main Arduino}\\
{\small Centrale Nantes}\\
\small P.-E. Hladik, \small{\texttt{pierre-emmanuel.hladik@ec-nantes.fr}}\\
---\\
{\scriptsize  Version \version (\today)}\\
}


\begin{document}

\maketitle

%\begin{table}[htdp]
%\caption{Suivi des modifications}
%\begin{center}
%\begin{tabular}{|c|c|c|p{8cm}|}
%\hline
%Date & Version & Auteur & Modifications\\
%\hline
%JJ/MM/12 & 0.3.X & P.-E. Hladik&  \\
%\hline
%\end{tabular}
%\end{center}
%\label{default}
%\end{table}%

%%%%%%%%%%%%%%%%%%%%%%%%%
  \section{Objectifs pédagogiques}
%%%%%%%%%%%%%%%%%%%%%%%%%

\begin{itemize}
	\item utiliser l'IDE Arduino
	\item utiliser un analyseur logique et PulseWiew
	\item bonne pratique (simple) pour implémenter des fonctions périodiques
\end{itemize}

%%%%%%%%%%%%%%%%%%%%%%%%%
  \section{Programmer sur Arduino}
%%%%%%%%%%%%%%%%%%%%%%%%%

Faire clignoter sur la Teensy une led à $2~Hz$ et une autre de manière aléatoire avec un délai entre 200 et 700 ms (code initial du TD 0 sur hippocampus ). Le code disponible (ci-dessous) est uniquement fourni pour vous aider à prendre en main les fonctions existantes pour contrôler les leds.

\begin{lstlisting}[language=C]
#define PIN_LED0 3
#define PIN_LED1 4

void setup() {
  pinMode(PIN_LED0, OUTPUT); // Initialisation de la broche 3
  pinMode(PIN_LED1, OUTPUT); // Initialisation de la broche 4 
  randomSeed(analogRead(0)) ; // Initialisation du generateur aleatoire
}

void loop() {
  digitalWrite (PIN_LED1, 1) ; // mise a l'etat haut de la broche 3
  delay(500); // attente de 500 ms
  digitalWrite (PIN_LED0, 0) ; // mise a l'etat bas de la broche 4
  delay(200+random(500)); // attente aleatoire de 200 a 699 ms
}
\end{lstlisting}
  
La fonction \href{https://www.arduino.cc/reference/en/language/functions/time/delay/}{\tt delay} bloque l'exécution, regarder plutôt \href{https://www.arduino.cc/reference/en/language/functions/time/millis/}{\texttt{millis}}.

%%%%%%%%%%%%%%%%%%%%%%%%%
  \section{Utilisation de PulseView}
%%%%%%%%%%%%%%%%%%%%%%%%%

Utiliser l'analyseur logique pour visualiser les signaux (voir le guide dédié à PulseView sur hippocampus). Vous pouvez les observer directement sur les fiches bananes à l'arrière de la carte. 


\end{document}